\chapter{Feasibility Spike}

\section{Purpose}

The feasibility spike validates the riskiest architectural assumption of this project: that
nopCommerce can remain the commerce core of the VerdeMart omnichannel scenario while a
durable, independently deployable integration layer is built around it without rewriting
checkout or the core order model.

The spike was conducted in two stages. The first stage, completed at the architecture
checkpoint, was a source-code inspection that confirmed the existence of the required seams
in the nopCommerce baseline. The second stage, completed for the final delivery, is a full
runtime implementation that exercises those seams under the mandatory Scenario~C pressure
points: warehouse unavailability, shipping outage, inventory staleness, and cross-channel
inventory conflict.

\section{What Was Tested}

The runtime spike implements and exercises the following capabilities end to end:

\begingroup
\small
\setlength{\tabcolsep}{4pt}
\renewcommand{\arraystretch}{1.15}
\begin{center}
\begin{tabular}{|p{0.28\textwidth}|p{0.35\textwidth}|p{0.27\textwidth}|}
\hline
\textbf{Spike question} & \textbf{How tested at runtime} & \textbf{Result} \\
\hline
Does order creation atomically produce an outbox record without blocking checkout? &
Real checkout via the nopCommerce one-page checkout. Outbox record verified in DB
immediately after order confirmation. &
Yes. \texttt{OrderProcessingService} writes the order and the outbox record in the same
transaction. Checkout response is returned before the warehouse is contacted. \\
\hline
Does the Integration Worker dispatch fulfillment commands asynchronously via RabbitMQ? &
Worker logs inspected after order placement. RabbitMQ management UI queues observed. &
Yes. Worker picks up the outbox record within polling interval, publishes
\texttt{FulfillmentRequested} to \texttt{fulfillment.requests}, and marks the record
\texttt{Published}. \\
\hline
Does the retry policy and circuit breaker protect the system under warehouse failure? &
Warehouse stub set to \texttt{failed} mode. Orders placed. Worker logs and Operations View
observed until circuit breaker opened. &
Yes. Exponential backoff retries are scheduled; after the configured failure threshold the
circuit breaker opens; subsequent calls are skipped; the Operations View shows the adapter
state and retry count. \\
\hline
Does the system recover automatically when the warehouse becomes available? &
Warehouse stub restored to \texttt{normal}. Cooldown elapsed. Worker probe observed in logs. &
Yes. Circuit breaker enters \texttt{HalfOpen} state, sends a probe, succeeds, and closes.
Pending outbox records are dispatched automatically. \\
\hline
Does the dead-letter queue capture messages that exhaust all retries? &
Shipping stub set to \texttt{outage} mode. Order placed. Worker left to exhaust max retry
attempts. &
Yes. After max attempts a \texttt{DeadLetterRecord} is created with adapter, failure reason,
and correlation ID. The Operations View surfaces the entry for operator action. \\
\hline
Does the requeue action recover a dead-letter record without operator scripting? &
Requeue button clicked in Operations View. Shipping stub restored. &
Yes. A new outbox record is created from the dead-letter payload. The worker dispatches it
on the next cycle and marks it \texttt{Published}. \\
\hline
Is inventory staleness detected and surfaced within the configured threshold? &
Inventory stub set to \texttt{unavailable}. Staleness threshold set to 30~s for the spike. &
Yes. After 30~s without a confirmed stock update, \texttt{IsStale} is set on the affected
\texttt{InventoryProjectionRecord} rows. The Operations View Inventory tab shows yellow
rows. \\
\hline
Is a POS--WMS inventory conflict detected and the checkout quantity bounded? &
POS stock injected via \texttt{POST /stock/pos-report} with a quantity diverging from the
WMS value by more than the configured tolerance. &
Yes. The worker detects the divergence on the next sync cycle, sets \texttt{ConflictFlag},
records \texttt{PendingReconciliation}, and writes the lower quantity to
\texttt{Product.StockQuantity} to prevent overselling. \\
\hline
\end{tabular}
\end{center}
\endgroup

\section{What Worked}

All eight runtime questions were answered positively. The key findings are:

\textbf{Checkout is fully decoupled from operational systems.} A customer order is accepted
and the HTTP response returned before the warehouse, shipping provider, or inventory system
is contacted. The outbox record written atomically with the order is the only coupling point.
This directly satisfies QAS~1 and QAS~3: warehouse unavailability cannot block checkout, and
the customer-facing path latency is unaffected by backend synchronisation.

\textbf{The retry and circuit breaker model behaves as designed.} Under warehouse failure,
the worker applies exponential backoff with the configured initial delay and multiplier. After
the failure threshold is reached, the circuit opens and no further calls are made until the
cooldown expires. After recovery, the half-open probe succeeds and the circuit closes without
operator intervention. This end-to-end lifecycle was observed in runtime logs and confirmed
in the Operations View.

\textbf{The dead-letter and requeue flow is operator-accessible.} When a message exhausts
all retries, the Operations View surfaces it with a one-click requeue action. The requeued
message is treated as a new outbox record and dispatched normally. No database access or
scripting is required from the operator.

\textbf{Inventory freshness and conflict detection work at configured thresholds.} Staleness
is detected within the configured 30-second threshold in the demo environment. Cross-channel
conflicts are detected within one sync cycle (10~seconds in the demo environment) and the
conservative lower quantity is immediately applied to the checkout stock gate.

\textbf{The nopCommerce seams identified in the checkpoint inspection are correct.} The
\texttt{PlaceOrderAsync} method in \texttt{OrderProcessingService} provides the atomic
write seam. The \texttt{ScheduleTaskRunner} locking mechanism is not used by the Integration
Worker (which runs as an independent process), but the seam inspection confirmed that the
nopCommerce scheduled-task infrastructure is available as a fallback execution host if needed.
The gap in the in-process \texttt{EventPublisher} identified at the checkpoint is correctly
addressed by the outbox: no external integration depends on the in-process event bus.

\section{What Did Not Work Yet}

The following items were identified as out of scope for the spike and are recorded as known
limitations in \texttt{docs/evidence/known-limitations.md}:

\begin{itemize}
    \item The POS source has no autonomous event emitter. Conflict detection is triggered
    via \texttt{POST /stock/pos-report} on the inventory stub, which simulates a POS stock
    report without a full POS lifecycle.
    \item Circuit breaker state is not persisted across Integration Worker restarts. On
    restart the registry initialises to \texttt{Closed} regardless of the last recorded
    state.
    \item Idempotency records are not cleaned up by a scheduled task. The \texttt{ExpiresAtUtc}
    column is populated but no cleanup job runs in the spike environment.
    \item The \texttt{IsStale} and \texttt{ConflictFlag} states are surfaced in the Operations
    View but are not exposed on the customer-facing product pages. Customers see the
    conservative stock quantity but not an explicit freshness indicator.
    \item Performance measurements under sustained load (QAS~3 response measures) have not
    been collected. The decoupling is structural and visible in the code; the numeric
    validation requires a dedicated load-testing environment.
\end{itemize}

\section{Evidence}

Runtime evidence is collected in \texttt{docs/evidence/} and covers the following scenarios:

\begingroup
\small
\setlength{\tabcolsep}{4pt}
\renewcommand{\arraystretch}{1.15}
\begin{center}
\begin{tabular}{|p{0.30\textwidth}|p{0.30\textwidth}|p{0.30\textwidth}|}
\hline
\textbf{Scenario} & \textbf{Evidence file} & \textbf{What it shows} \\
\hline
Normal flow & \texttt{screenshots/05-normal-flow.png} &
Outbox records with \texttt{Published} status; all circuit breakers \texttt{Closed}. \\
\hline
RabbitMQ queues & \texttt{screenshots/06-rabbitmq-queues.png} &
\texttt{fulfillment.requests}, \texttt{shipping.requests}, \texttt{fulfillment.dead}, and
\texttt{shipping.dead} queues present and active. \\
\hline
Warehouse failure & \texttt{screenshots/07-warehouse-failure-cb-open.png} &
Outbox records in \texttt{Retrying} and \texttt{Failed} states; warehouse circuit breaker
in \texttt{Open} state. \\
\hline
Warehouse recovery & \texttt{screenshots/08-warehouse-recovery.png} &
All circuit breakers returned to \texttt{Closed}; pending records dispatched. \\
\hline
Dead-letter creation & \texttt{screenshots/09-deadletter-created.png} &
Dead letter record with \texttt{EscalationState = New} after shipping outage exhausts
retries. \\
\hline
Dead-letter requeue & \texttt{screenshots/10-deadletter-requeued.png} &
Dead letter \texttt{EscalationState = Requeued}; new outbox record \texttt{Published}. \\
\hline
Inventory staleness & \texttt{screenshots/03-inventory-stale.png} &
\texttt{IsStale = true} on inventory projection rows after threshold exceeded. \\
\hline
Inventory conflict & \texttt{screenshots/04-inventory-conflict.png} &
\texttt{ConflictFlag = true} on affected rows; checkout quantity bounded to lower value. \\
\hline
Worker logs & \texttt{logs/warehouse-failure.txt} &
Retry attempts with exponential delay, circuit breaker open and close events. \\
\hline
\end{tabular}
\end{center}
\endgroup

The complete runtime environment is reproducible from the repository root using
\texttt{make up} followed by \texttt{./install-nopcommerce.sh}. All pressure scenarios are
exercisable via \texttt{make warehouse-fail}, \texttt{make shipping-fail},
\texttt{make stale-start}, and \texttt{make conflict-inject} as documented in
\texttt{docs/evidence/demo-script.md}.
